\section{METHODOLOGY}

\subsection{Study Design}
CurricuPath was tested on 139 synthetic cases for the Wits EFA03 and EFA04
programmes, based on handbook rules and School registration workbooks. No
real student records were used. The case file was frozen by hash before the
reported runs. The cases cover programme membership, prerequisites,
corequisites, timetable clashes, missing student information, incomplete
evidence and institutional discretion. Each case was phrased with course
codes, module names and informal language. The module-name and informal
levels use only three and four sentence templates, so they test coverage of
a few phrasings rather than ambiguity in general.

CurricuPath was compared with a version without clarification, a rules-only
evaluator, a Handbook RAG baseline and a structured-rules LLM baseline.
CurricuPath and the first two alternatives were tested on all 139 cases.
API costs and quotas limited the two LLM baselines to a 40-case subset,
drawn with a fixed seed before any baseline result was read, preserving the
outcome proportions: 14 eligible, 20 ineligible and 6 manual review
required. All five systems were compared on this subset.

\subsection{Case Construction and Controls}
Cases were constructed from the 2026 EBE handbook \cite{wits2026ebe},
processed curriculum rules and School of EIE registration workbooks. Cases,
student profiles, questions and reference outcomes were frozen before the
reported runs, and all systems received the same profile facts, including
declared review reasons. The systems were not frozen. After the first
results, the parser was extended to recognise registration wording (for
example, ``allowed to register for''), the School grid prerequisites were
loaded, and review reasons were passed to the engine. Scores after these
changes are development-informed; the earlier scores (Macro-F1 0.00 at the
module-name level and 0.35 at the informal level) are retained. The
workbooks, supplied by the School's Admin Manager on 30 September 2026, were
the primary reference. As no independent academic reviewer was obtained,
the results measure agreement with this School reference rather than
independently validated academic accuracy.

\subsection{System Comparison}
CurricuPath applies fixed eligibility rules and asks a follow-up question when
a material student fact is missing. In the reported runs it read questions
with a deterministic parser and called Gemini only for messages the parser
could not resolve. Its no-clarification version uses the same rules but never
asks. The rules-only evaluator applies the rules directly to the structured
profile. Handbook RAG answers from retrieved handbook passages, while the
structured-rules LLM receives the rule table and student profile.

Both baselines used \texttt{gemini-\allowbreak 3.5-\allowbreak flash-\allowbreak lite}
\cite{googleGemini35FlashLite}. The planned model, \texttt{gemini-3.6-flash},
could not be used within the free-tier quota, so the same lower-cost model
served both baselines; their results therefore describe a low-cost model.
Temperature zero and one run per case were used. The initial baseline runs
lacked the School grid prerequisites and corequisites, so both were rerun with
that evidence added. The Handbook RAG arm is therefore no longer
handbook-only. Both reruns completed; one structured-rules run (informal
wording) failed with a network disconnect and counts as an invalid output.

\subsection{Reference Decisions}
Each case has one reference outcome: \textit{eligible}, \textit{ineligible} or
\textit{manual review required}, taken from the School registration
workbooks. Handbook-only labels are stored with each case but were not
analysed. No reviewer agreement is reported because no independent reviewers
were available. A timetable clash involving a module the student was
repeating made the case ineligible; the clash rule is a reviewer statement
still awaiting written confirmation.

\subsection{Evaluation Measures}
The main measure was Macro-F1, which weights the three outcomes equally
\cite{sokolova2009performance}, with bootstrap intervals over cases. We also
report balanced accuracy, precision and recall by outcome, confusion matrices
and raw counts. The unsafe false-positive rate is the proportion of
\textit{ineligible} or \textit{manual review required} cases labelled
\textit{eligible}. Results are also split by whether the handbook alone could
answer the case and by wording level, and abstention is measured as precision
and recall of deferring to manual review, with over- and under-referral
rates. Proportions use Wilson intervals. Invalid outputs remain in the totals.
Response time and provider failures were recorded.

Clarification was tested separately, because the main set's clarification
cases use a scripted reply the system cannot read. Two dedicated sets of 72
cases each contain 33 cases missing one needed fact, 33 matched controls with
the fact known, and 6 with an unrelated unknown. A scripted reply supplies the
fact and the final decision is scored. In the first set, 8 of 33 cases ended
undecided because of an extra clash-partner question; the second removed that
side effect. Their labels are developer-proposed, pending academic review.

\subsection{Interpretation Constraints}
The rules-only evaluator was built from the same School workbooks used to set
the reference outcomes. Its perfect agreement shows faithful implementation of
those workbooks, not independently established correctness, and CurricuPath's
scores share this limit. The second baseline runs have a similar limitation:
the cases had already been inspected, and the added grid evidence came from
those workbooks. The baselines used one low-cost model, one run per case and
40 cases, six of them manual review, so their intervals are wide. Stronger
claims require new cases and reference outcomes established independently of
the systems.
